\documentclass{fisatprojectfinal}
\renewcommand{\baselinestretch}{1.25}
\usepackage{setspace}
\usepackage[utf8]{inputenc}
\usepackage{xcolor}
\usepackage[a4paper,left=3.8cm,right=2.2cm,top=2.8cm,bottom=2.8cm,bindingoffset=0.7cm]{geometry}
\usepackage{microtype}
\usepackage{listings}
\usepackage{caption}

\setlength{\emergencystretch}{3em}
\pretolerance=1000
\tolerance=2000

\lstset{
	basicstyle=\ttfamily\small,
	breaklines=true,
	breakatwhitespace=true,
	columns=fullflexible,
	frame=single,
	keepspaces=true,
	showstringspaces=false,
	tabsize=4
}

\title{College Canteen Face Detection and Visit Logging System}
\team{Team Member 1 \\ Team Member 2 \\ Team Member 3 \\ Team Member 4}
\author{Team Member 1}
\regno{FITCS20001}
\bibliographystyle{plain}

\begin{document}
\maketitle
\makecert

\newpage
\pagenumbering{roman}
\setcounter{page}{1}
\newgeometry{top=4cm,bottom=0.1cm}
\thispagestyle{empty}
\renewcommand\abstractname{ABSTRACT}
\begin{abstract}
	\vspace{2cm}
	\begin{onehalfspace}
		This mini project presents a real-time College Canteen Face Detection and Visit Logging System that addresses recurrent issues related to false canteen usage complaints, missing traceability of visits, and manual monitoring overhead. The application combines computer vision, facial embedding-based recognition, and database-backed logging to provide an automated mechanism for recording who entered the canteen and when.

		The primary objective is to build a practical campus-level solution that can detect faces from a live camera or RTSP stream, identify known students, label unknown visitors, and persist visit events with timestamps for reporting and review. The project uses OpenCV for video input and preprocessing, DeepFace with a FaceNet-based embedding model for recognition, and SQLite for lightweight persistent storage. The system uses a desktop interface built with Tkinter for operation and monitoring.

		The implemented pipeline includes student registration (from live capture and uploaded images), frame-wise detection with performance optimizations, cosine-similarity matching with margin constraints, visit deduplication using time windows, and statistical summaries such as total visits, unique visitors, and unknown detections. RTSP support and auto-reconnect logic make the solution deployable in practical network-camera setups.

		Testing was carried out through module-level and workflow-level scenarios covering registration, detection accuracy under varying conditions, logging, filtering, export, and stream stability. The developed system demonstrates that a low-cost and modular architecture can provide operationally useful accountability and analytics for institutional canteen environments.
	\end{onehalfspace}
\end{abstract}

\newpage
\renewcommand\abstractname{Contribution by Author}
\thispagestyle{plain}
\begin{abstract}
	\vspace{5cm}
	Author 1 Contribution: Core architecture, face recognition pipeline integration, and system testing coordination. \\
	Author 2 Contribution: Database schema design, logging/reporting modules, and CSV export features. \\
	Author 3 Contribution: Tkinter GUI development, user interaction flow, and student/log management screens. \\
	Author 4 Contribution: RTSP integration, performance tuning, deployment scripts, and documentation. \\
	\vspace{1cm}
	\begin{flushright}
		Team Member 1 \\
		Team Member 2 \\
		Team Member 3 \\
		Team Member 4
	\end{flushright}
\end{abstract}

\newpage
\renewcommand\abstractname{ACKNOWLEDGMENT}
\thispagestyle{plain}
\begin{abstract}
	\vspace{5cm}
	We express our sincere gratitude to our project guide, faculty members, and department for their continuous support in completing this mini project. Their feedback on problem formulation, implementation approach, and testing strategy helped us improve both the technical quality and the practical relevance of the work. We also thank our classmates and lab staff for assisting with testing environments, camera setup, and validation trials.

	We acknowledge the open-source communities behind Python, OpenCV, DeepFace, and SQLite, which provided dependable tools and documentation for this work.
	\vspace{1cm}
	\begin{flushright}
		Team Members
	\end{flushright}
\end{abstract}

\newpage
\thispagestyle{empty}
\restoregeometry
\newpage

\thispagestyle{empty}
\setcounter{page}{0}
\tableofcontents
\thispagestyle{empty}
\newpage

\cleardoublepage
\addcontentsline{toc}{chapter}{\listfigurename}
\thispagestyle{empty}
\pagenumbering{gobble}
\listoffigures
\thispagestyle{empty}
\newpage

\cleardoublepage
\addcontentsline{toc}{chapter}{\listtablename}
\thispagestyle{empty}
\pagenumbering{gobble}
\listoftables
\thispagestyle{empty}
\newpage

\chapter{Introduction}
\pagenumbering{arabic}
\setcounter{page}{1}
\renewcommand{\baselinestretch}{1.50}

\section{Overview}
The project focuses on building a reliable and real-time face-based visit tracking system for a college canteen. In many institutions, canteen disputes occur when students claim they did not visit the canteen on a given day. Manual records are incomplete and surveillance footage is difficult to audit quickly. This creates an accountability and transparency gap.

The proposed system captures video frames from either a local webcam or an RTSP stream, detects faces, and recognizes registered students by comparing facial embeddings. Each successful detection can trigger a visit log entry with metadata such as date, entry time, and optional screenshot path.

The starting point of the project was the observation that existing canteen operations lacked:
\begin{itemize}
	\item a structured identity-linked visit log,
	\item automated handling of unknown visitors,
	\item periodic statistics and exports for administrative review.
\end{itemize}

The implemented work includes a complete end-to-end solution with registration workflows, real-time recognition, logging, GUI dashboards, and reporting utilities.

Social relevance comes from improving fairness in complaint handling and reducing manual workload for canteen and college administrators. Applications include canteen analytics, basic security monitoring, and research demonstrations in computer vision.

Ethical and professional considerations addressed in this project include controlled storage of face data, consent-driven registration, role-based access to logs, and recommendations for retention policies.

\begin{figure}[h!]
	\begin{center}
		\fbox{\rule{0pt}{2in}\rule{0.8\linewidth}{0pt}}
		\caption{High-level workflow: capture, detect, recognize, log, and report}
	\end{center}
\end{figure}

\section{Problem Statement}
College canteens often face operational complaints related to billing disputes, crowd management, and unverified claims about student presence. In typical workflows, there is no structured and searchable record that maps identity to visit time in a non-intrusive and automated manner. Existing CCTV systems are passive and require manual review, which is time-consuming and error-prone.

The general problem is the absence of an intelligent, low-cost, and deployable mechanism for accountability in student service spaces. This affects administrators, canteen staff, and students due to delayed dispute resolution and limited data for management decisions.

The specific problem addressed in this project is to design and implement a canteen-focused face recognition system that can:
\begin{itemize}
	\item identify registered students in real time,
	\item flag unknown faces,
	\item store visit records with timestamps,
	\item provide daily statistics and filtered reports,
	\item support both webcam and RTSP-based deployments.
\end{itemize}

\section{Objectives}
The major objectives of this mini project are listed below:
\begin{itemize}
	\item Develop a real-time face detection and recognition pipeline suitable for college canteen monitoring.
	\item Implement student registration and embedding storage using a structured database model.
	\item Build visit logging with duplicate suppression using configurable time intervals.
	\item Support camera sources including local webcam and RTSP network streams.
	\item Provide intuitive interfaces for operation, monitoring, filtering, and export.
	\item Generate usable statistics for daily review and management decisions.
	\item Design the system to be modular, maintainable, and extensible for future upgrades.
\end{itemize}

\section{Scope of the project}
The current scope includes face-based entry event tracking in a single canteen environment. It covers registration, recognition, logging, analytics, and report export at prototype level. The solution is targeted for educational and controlled campus deployment.

Out-of-scope items for this phase include biometric liveness detection, multi-campus distributed synchronization, legal-grade evidence handling, and payment integration.

\section{Social Relevance}
This project contributes to a more transparent and accountable campus service ecosystem. It supports fair handling of disputes, improves trust between stakeholders, and reduces manual administrative burden. It also offers learning opportunities for students in AI-enabled public utility systems.

\section{Organization of the report}
This mini project report is organized into seven chapters. Chapter 1 introduces the problem, objectives, scope, and social relevance. Chapter 2 reviews related literature and positions the proposed system. Chapter 3 presents requirements, architecture, design documents, APIs, and algorithms. Chapter 4 explains implementation details including modules, tools, and deployment. Chapter 5 covers test planning, scenarios, and results. Chapter 6 summarizes outputs, comparative analysis, and sustainability impact. Chapter 7 concludes the work and highlights future scope.

\chapter{Literature Review}
After introducing the project context, this chapter discusses existing methods and technologies relevant to face detection, embedding-based recognition, and practical deployment.

\section{Related work}
DeepFace introduced high-performance deep architectures for face verification and demonstrated near-human-level performance on benchmark datasets \cite{deepface_framework}. FaceNet proposed a unified embedding space where Euclidean distance can directly represent face similarity \cite{facenet_paper}. YOLO-based object detection methods established real-time detection capabilities suitable for live vision systems \cite{yolo_paper}.

For implementation-level development, OpenCV provides robust APIs for image processing, camera handling, and DNN inference \cite{opencv_library}. SQLite offers lightweight relational persistence suitable for local desktop deployments \cite{sqlite_doc}.

\section{Comparison of related works}
\begin{table}[h!]
\centering
\caption{Comparison of major approaches relevant to this project}
\begin{tabular}{|l|l|l|l|}
\hline
Approach & Strengths & Limitations & Relevance \\
\hline
Traditional CCTV review & Always available video & Manual, non-searchable & Low \\
\hline
Classical Haar-only system & Fast on CPU & Lower recognition robustness & Medium \\
\hline
Deep embedding systems & Better identity matching & Compute-intensive & High \\
\hline
YOLO-based detection stack & Real-time performance & Needs tuning for devices & High \\
\hline
\end{tabular}
\end{table}

\section{Proposed System}
The proposed system combines real-time frame acquisition, face detection, embedding extraction, identity matching, and event logging in a modular architecture.

Methodology summary:
\begin{itemize}
	\item Capture frames from webcam or RTSP source.
	\item Detect candidate faces using DNN/Haar-based strategies.
	\item Enhance cropped faces and compute embeddings using DeepFace with FaceNet model.
	\item Match against known embeddings using cosine similarity and margin filtering.
	\item Log known or unknown events to database with date/time metadata.
	\item Expose outputs through a desktop interface, including statistics and CSV exports.
\end{itemize}

\chapter{Design Methodologies}

\section{Software Requirement Specifications}
The software requirements for this project were derived from operational needs of canteen monitoring and practical constraints of college lab infrastructure.

\subsection{Functional Requirements}
\begin{itemize}
	\item Start and stop live detection from webcam or RTSP source.
	\item Register students with ID, name, department, year, face image, and embedding.
	\item Register from both live capture and uploaded image file.
	\item Detect and identify faces in near-real-time.
	\item Log visits with timestamps, identity labels, and optional screenshot path.
	\item Maintain unknown-face records for later review.
	\item Display live statistics (known, unknown, total, FPS).
	\item Filter logs by date and student ID and export logs/students to CSV.
\end{itemize}

\subsection{Non-Functional Requirements}
\begin{itemize}
	\item Performance: interactive processing with frame skipping and threaded recognition.
	\item Reliability: reconnect behavior for unstable RTSP sources.
	\item Usability: simple desktop interface for non-expert users.
	\item Maintainability: modular Python files for database, recognition, UI, and utilities.
	\item Portability: executable on standard Windows systems with Python environment.
	\item Data integrity: structured storage using relational schema with key constraints.
\end{itemize}

\section{Software Design Document}

\subsection {System Features}
The major system features are:
\begin{itemize}
	\item Real-time detection and recognition.
	\item Student lifecycle management (add, update, delete).
	\item Visit and unknown-face logging.
	\item Daily analytics and report export.
	\item RTSP support with reconnect and buffering controls.
\end{itemize}

\subsection{System Architecture Design}
The system follows a layered pipeline architecture:
\begin{enumerate}
	\item Input Layer: Webcam/RTSP stream acquisition.
	\item Vision Layer: Face detection and quality enhancement.
	\item Recognition Layer: Embedding extraction and matching.
	\item Persistence Layer: SQLite storage for students and logs.
	\item Presentation Layer: Tkinter GUI.
\end{enumerate}

\begin{figure}[h!]
	\begin{center}
		\fbox{\rule{0pt}{2in}\rule{0.8\linewidth}{0pt}}
		\caption{System architecture showing data flow across layers}
	\end{center}
\end{figure}

\subsection{Constraints}
\begin{itemize}
	\item Recognition quality depends on lighting, pose, and registration image quality.
	\item Low-end hardware may reduce frame-rate under heavy load.
	\item Network variability affects RTSP stability and latency.
	\item Privacy and consent requirements constrain deployment in real environments.
\end{itemize}

\subsection{Application Architecture Design}
Application modules are separated by concern:
\begin{itemize}
	\item \textbf{config.py}: runtime parameters and thresholds.
	\item \textbf{database.py}: schema setup and CRUD operations.
	\item \textbf{face\_recognition\_module.py}: detection, embedding, matching pipeline.
	\item \textbf{face\_matcher.py}: similarity and best-match logic.
	\item \textbf{gui\_app.py}: desktop management interface.
	\item \textbf{utils.py}: report generation and export utilities.
\end{itemize}

\subsection {API Design}
Module-level interfaces used in this desktop implementation include:
\begin{itemize}
	\item \texttt{main.py} detection controls: start camera, stop camera, run detection loop.
	\item \texttt{gui\_app.py} actions: register student, refresh students, filter logs, export CSV.
	\item \texttt{database.py} functions: add student, update student, delete student, log visit.
	\item \texttt{utils.py} functions: generate report, export logs, export students.
\end{itemize}

\subsection {Database Design}
The project uses SQLite with three primary tables:
\begin{itemize}
	\item \textbf{students}: student profile and embedding storage.
	\item \textbf{visit\_logs}: entry/exit events and status labels.
	\item \textbf{unknown\_faces}: unregistered face snapshots and embeddings.
\end{itemize}

\begin{table}[h!]
\centering
\caption{Students table design}
\begin{tabular}{|l|l|l|}
\hline
Field & Type & Description \\
\hline
id & INTEGER & Primary key \\
student\_id & TEXT & Unique student identifier \\
name & TEXT & Student name \\
department & TEXT & Department name \\
year & INTEGER & Year of study \\
face\_embedding & TEXT & JSON encoded embedding vector \\
face\_image\_path & TEXT & Stored image path \\
created\_at & TIMESTAMP & Record creation time \\
updated\_at & TIMESTAMP & Last update time \\
\hline
\end{tabular}
\end{table}

\begin{table}[h!]
\centering
\caption{Visit logs table design}
\begin{tabular}{|l|l|l|}
\hline
Field & Type & Description \\
\hline
id & INTEGER & Primary key \\
student\_db\_id & INTEGER & Foreign key to students \\
student\_id & TEXT & Student identifier \\
student\_name & TEXT & Detected name \\
entry\_time & TIMESTAMP & Visit start time \\
exit\_time & TIMESTAMP & Visit end time \\
duration\_minutes & INTEGER & Visit duration \\
date & DATE & Local date \\
is\_known & INTEGER & Known/unknown flag \\
screenshot\_path & TEXT & Optional frame snapshot path \\
\hline
\end{tabular}
\end{table}

\subsection {Technology Stack}
\begin{itemize}
	\item Language: Python 3.10+
	\item Computer Vision: OpenCV
	\item Recognition: DeepFace (FaceNet embedding model)
	\item Detection Model Support: YOLOv8 assets and DNN detector files
	\item Desktop UI: Tkinter
	\item Database: SQLite
	\item Numeric Processing: NumPy
\end{itemize}

\section{Use case Diagrams/Data flow Diagrams}
\textbf{Use Case Summary:}
\begin{itemize}
	\item Admin registers student profiles.
	\item System captures and processes live frames.
	\item System logs known and unknown visits.
	\item Admin reviews logs and exports reports.
\end{itemize}

\textbf{Data Flow Summary:}
\begin{enumerate}
	\item Video frame input from camera source.
	\item Face region extraction and preprocessing.
	\item Embedding generation and matching with known vectors.
	\item Log event generation and database update.
	\item Dashboard retrieval and report export.
\end{enumerate}

\begin{figure}[h!]
	\begin{center}
		\fbox{\rule{0pt}{2in}\rule{0.8\linewidth}{0pt}}
		\caption{Placeholder for use case and data-flow diagrams}
	\end{center}
\end{figure}

\subsection*{Detailed Use Case Description}
\textbf{Actor: Administrator}
\begin{itemize}
	\item Pre-condition: System is installed and database is initialized.
	\item Main flow: Administrator launches the GUI, opens student tab, registers student identity, and validates successful insertion.
	\item Alternate flow: If no face is detected in registration frame, system prompts recapture and rejects incomplete records.
	\item Post-condition: Student metadata and embedding are stored and available for recognition.
\end{itemize}

\textbf{Actor: Detection Operator}
\begin{itemize}
	\item Pre-condition: Camera or RTSP source is active.
	\item Main flow: Operator starts detection, monitors recognized names and confidence, and observes auto logging.
	\item Alternate flow: If stream fails, reconnect logic is triggered; operator is notified in terminal/UI status indicators.
	\item Post-condition: Visit logs are persisted with time and identity/unknown state.
\end{itemize}

\textbf{Actor: Report Reviewer}
\begin{itemize}
	\item Pre-condition: Visit logs exist in the selected date range.
	\item Main flow: Reviewer filters logs by date/student ID, verifies entries, and exports CSV.
	\item Alternate flow: If no records are found, system returns empty dataset without failure.
	\item Post-condition: Report artifacts are generated for academic/administrative verification.
\end{itemize}

\subsection*{Data Dictionary Snapshot}
\begin{table}[h!]
\centering
\caption{Data dictionary summary for major entities}
\begin{tabular}{|l|l|l|}
\hline
Entity & Key Fields & Remarks \\
\hline
Student & student\_id, name, embedding & Core identity registry \\
\hline
Visit Log & entry\_time, is\_known, screenshot\_path & Operational tracking \\
\hline
Unknown Face & first\_seen, times\_seen & Follow-up candidate \\
\hline
Statistics & total, unique, unknown & Daily summary values \\
\hline
\end{tabular}
\end{table}

\section{Algorithms}
\textbf{Algorithm 1: Real-time detection and recognition}
\begin{enumerate}
	\item Read latest frame from buffer.
	\item Detect candidate face regions.
	\item For each face region, enhance and compute embedding.
	\item Compare embedding with known vectors using cosine similarity.
	\item Assign identity if threshold and margin conditions are satisfied.
	\item Render annotations and update detection counters.
\end{enumerate}

\textbf{Sample Input:} live frame from webcam/RTSP stream.

\textbf{Sample Output:} annotated frame with label \texttt{Name (confidence)} and updated visit status.

\textbf{Algorithm 2: Visit logging with time-window control}
\begin{enumerate}
	\item On successful recognition, fetch most recent visit of the student.
	\item Compute elapsed time from previous entry.
	\item If elapsed time is less than minimum threshold, skip duplicate log.
	\item Else insert a new visit entry into \texttt{visit\_logs}.
\end{enumerate}

\textbf{Sample Input:} student\_id and current timestamp.

\textbf{Sample Output:} inserted log row ID or no-op decision.

\chapter{Implementation}
Software implementation in this project converts the proposed architecture into a working prototype with modular code and reusable components.

\section{Implementation Details}
\subsection{Code Development}
Code was developed as independent Python modules to separate concerns and reduce coupling. The recognition module handles compute-heavy operations while database and UI modules remain isolated for maintainability.

\subsubsection{Set Coding Standards}
\begin{itemize}
	\item Function-level docstrings for key modules.
	\item Clear module naming based on responsibility.
	\item Configurable constants centralized in \texttt{config.py}.
	\item Defensive checks for camera failures and invalid inputs.
\end{itemize}

\subsection{Source Code Control Setup}
The project repository is organized with dedicated source files, documentation guides, model assets, and templates. Version control was used to manage feature increments such as RTSP support, recognition improvements, and performance tuning.

\subsection{Dataset}
The project uses a custom face dataset generated during registration. Each student profile stores a representative face image and a corresponding embedding vector in database storage. Unknown faces are optionally stored for later review.

\section{Libraries/Applications}
Major libraries and tools used:
\begin{itemize}
	\item OpenCV for capture, preprocessing, and visualization.
	\item DeepFace for embedding extraction.
	\item NumPy for vector operations.
	\item Tkinter for desktop dashboard.
	\item SQLite for local persistence.
\end{itemize}

Applications used during testing and deployment:
\begin{itemize}
	\item OBS RTSP plugin for stream simulation and validation.
	\item Windows command shell / PowerShell for script execution.
\end{itemize}

\section{Deployment}
Deployment procedure:
\begin{enumerate}
	\item Install dependencies using \texttt{pip install -r requirements.txt}.
	\item Run \texttt{python setup.py} for first-time setup.
	\item Start desktop app with \texttt{python gui\_app.py} or console mode with \texttt{python main.py}.
	\item For network camera mode, enable RTSP in configuration and verify with \texttt{python test\_rtsp.py}.
\end{enumerate}

\chapter{Testing}

\section{Test plan}
\subsubsection*{Test objectives}
Validate functionality, stability, and usability of registration, recognition, logging, and reporting flows.

\subsubsection*{Test scope}
Testing included unit-level utilities, module integration, camera/RTSP handling, and UI workflows.

\subsubsection*{Test approach}
Combined manual scenario testing and script-based checks using provided test modules.

\subsubsection*{Test environment}
Windows system with Python environment, webcam and optional RTSP source, SQLite local storage.

\subsubsection*{Test cases}
\begin{itemize}
	\item Student registration from live frame.
	\item Student registration from uploaded image.
	\item Recognition of known student.
	\item Unknown-face detection and logging.
	\item Duplicate suppression within configured time gap.
	\item Log filtering and CSV export.
	\item RTSP reconnect behavior on stream interruption.
\end{itemize}

\subsubsection*{Detailed Test Case Matrix}
\begin{table}[h!]
\centering
\caption{Functional test case matrix}
\begin{tabular}{|p{1.2cm}|p{4.2cm}|p{4.0cm}|p{3.6cm}|}
\hline
TC ID & Scenario & Expected Result & Status \\
\hline
TC-01 & Register known student from webcam frame & Student added to DB with embedding and image path & Pass \\
\hline
TC-02 & Register known student from uploaded image & Student profile created/updated successfully & Pass \\
\hline
TC-03 & Detect registered student in live stream & Name and confidence shown, visit logged once in threshold window & Pass \\
\hline
TC-04 & Detect unknown face in live stream & Unknown state shown and unknown logging pathway triggered & Pass \\
\hline
TC-05 & Repeated detection within minimum interval & Duplicate suppression prevents excessive logs & Pass \\
\hline
TC-06 & Filter logs by date and ID & Correct subset returned in GUI table & Pass \\
\hline
TC-07 & Export logs to CSV & CSV file generated with expected columns & Pass \\
\hline
TC-08 & RTSP disconnection and reconnect & Stream reconnect attempts initiated and resumed if source returns & Pass \\
\hline
TC-09 & Startup without available camera & Error message shown, application remains stable & Pass \\
\hline
TC-10 & Delete student and refresh list & Student removed and model cache reloaded & Pass \\
\hline
\end{tabular}
\end{table}

\subsubsection*{Test execution}
Executed through \texttt{test\_system.py}, \texttt{test\_rtsp.py}, and interactive GUI sessions.

\subsubsection*{Test schedule}
Testing was carried out iteratively after each major module integration phase.

\subsubsection*{Risks and Contingencies}
\begin{itemize}
	\item Risk: poor lighting reduces face quality. Contingency: image enhancement and better placement.
	\item Risk: RTSP packet loss. Contingency: reconnect and transport configuration.
	\item Risk: false positives. Contingency: threshold tuning and margin checks.
\end{itemize}

\subsubsection*{Conclusion}
The test plan verified that the prototype is functionally complete for canteen visit tracking and can be further tuned for environment-specific deployments.

\section{Testing Scenarios}
Scenario-based validation included:
\begin{itemize}
	\item Normal operation with registered students.
	\item Mixed known and unknown crowd conditions.
	\item Stream disconnection and auto-reconnect behavior.
	\item High-load periods with frequent frame events.
\end{itemize}
Diagrams and screenshots can be inserted in this section during final documentation submission.

\subsection*{Scenario Observations and Analysis}
\textbf{Scenario 1: Single user at optimal lighting}\newline
Observed confidence values were consistently high, with stable bounding boxes and minimal frame drops. System latency remained within acceptable interactive limits.

\textbf{Scenario 2: Multiple persons in frame}\newline
Detection handled multiple faces but recognition was intentionally conservative through threshold and margin checks to minimize false-positive identity assignment.

\textbf{Scenario 3: Low-light corridor entry}\newline
Face enhancement pipeline improved visibility. However, confidence reduction was observed in extreme low-light conditions, indicating deployment should ensure basic illumination.

\textbf{Scenario 4: RTSP jitter and packet delay}\newline
Frame fetch timeout and reconnect logic prevented complete application lock. Short disruptions were recovered without manual restart in most attempts.

\textbf{Scenario 5: High-traffic interval simulation}\newline
During repeated entries, deduplication logic reduced redundant rows and preserved meaningful visit analytics.

\begin{figure}[h!]
	\begin{center}
		\fbox{\rule{0pt}{2.4in}\rule{0.85\linewidth}{0pt}}
		\caption{Screenshot Placeholder 1: Live detection screen}
	\end{center}
\end{figure}

\begin{figure}[h!]
	\begin{center}
		\fbox{\rule{0pt}{2.4in}\rule{0.85\linewidth}{0pt}}
		\caption{Screenshot Placeholder 2: Student registration workflow}
	\end{center}
\end{figure}

\section{Unit Testing}
Unit checks focused on utility methods such as similarity computation, report generation, CSV export formatting, and database CRUD operations.

\begin{figure}[h!]
	\begin{center}
		\fbox{\rule{0pt}{2.4in}\rule{0.85\linewidth}{0pt}}
		\caption{Screenshot Placeholder 3: Unit test output}
	\end{center}
\end{figure}

\section{Integral Testing}
Integration testing validated end-to-end paths from frame capture through recognition and logging to UI update.

\begin{figure}[h!]
	\begin{center}
		\fbox{\rule{0pt}{2.4in}\rule{0.85\linewidth}{0pt}}
		\caption{Screenshot Placeholder 4: Integration test evidence}
	\end{center}
\end{figure}

\section{Functional Testing}
Functional tests confirmed that all user-facing operations are available and produce expected results across the desktop interface.

\begin{figure}[h!]
	\begin{center}
		\fbox{\rule{0pt}{2.4in}\rule{0.85\linewidth}{0pt}}
		\caption{Screenshot Placeholder 5: Functional test dashboard}
	\end{center}
\end{figure}

\section{Load Testing}
Load-oriented trials were performed by running prolonged streams and repeated recognitions to observe throughput, responsiveness, and logging consistency.

\subsection*{Load Test Summary Table}
\begin{table}[h!]
\centering
\caption{Indicative load testing summary in lab environment}
\begin{tabular}{|l|l|l|l|}
\hline
Test Window & Input Mode & Observation & Outcome \\
\hline
10 min & Webcam 640x480 & Stable UI and logging, no crash & Acceptable \\
\hline
20 min & RTSP source & Intermittent dropped frames, recovered by reconnect logic & Acceptable \\
\hline
30 min & Mixed user entries & Log persistence remained consistent & Acceptable \\
\hline
45 min & Continuous runtime & Memory usage increased moderately but remained manageable & Acceptable \\
\hline
\end{tabular}
\end{table}

\subsection*{Performance Notes}
\begin{itemize}
	\item Frame buffering with small queue size prevented stale frame accumulation.
	\item Threaded recognition reduced UI blocking and improved interaction smoothness.
	\item GPU-enabled paths provide noticeable throughput benefits where available.
	\item FPS values vary with camera quality, lighting, and active face count.
\end{itemize}

\begin{figure}[h!]
	\begin{center}
		\fbox{\rule{0pt}{2.4in}\rule{0.85\linewidth}{0pt}}
		\caption{Screenshot Placeholder 6: Load test monitoring}
	\end{center}
\end{figure}

\chapter {Result}

\section{Result summary}
The developed system successfully demonstrates real-time face-based canteen visit tracking with identity-linked logs and administrative visibility. Key achieved outcomes include:
\begin{itemize}
	\item practical registration-to-recognition workflow,
	\item persistent and queryable visit records,
	\item unknown visitor tracking,
	\item desktop-based operations,
	\item export-ready reporting utilities.
\end{itemize}

\subsection*{Quantitative Result Narrative}
Across repeated trials, the system maintained stable detection-to-logging flow in typical campus-like conditions. The log database remained queryable and consistent for daily summaries. Recognition stability improved after threshold tuning and enhanced registration quality. RTSP support enabled practical deployment alternatives beyond local webcams.

\subsection*{Qualitative Outcomes}
\begin{itemize}
	\item Better traceability for dispute handling compared with unstructured manual records.
	\item Easier operational oversight with integrated statistics and exports.
	\item Improved academic value through integration of CV, DB, UI, and deployment concerns.
\end{itemize}

\section {Comparison with existing system}
\begin{table}[h!]
\centering
\caption{Comparison with conventional manual monitoring}
\begin{tabular}{|l|l|l|}
\hline
Parameter & Existing Manual Approach & Proposed System \\
\hline
Identity traceability & Low & High \\
\hline
Record retrieval speed & Slow & Fast \\
\hline
Unknown visitor flagging & Not systematic & Built-in \\
\hline
Report generation & Manual & Automatic/CSV export \\
\hline
Scalability for daily review & Limited & Better \\
\hline
\end{tabular}
\end{table}

\subsection*{Comparative Discussion}
Manual systems depend heavily on human recollection and ad-hoc verification, which can introduce delay and inconsistency. The proposed implementation provides event-driven logs, faster retrieval, and clear filtering capabilities. While advanced biometric safeguards are outside this phase scope, current architecture already improves institutional process clarity.

\section{Sustainability Impact Assessment}
This project contributes to Sustainable Development Goal aligned digital campus governance by improving service accountability and reducing operational inefficiency.

\subsection{Contribution of Project Features to SDGs}
Selected SDG alignment: SDG 9 (Industry, Innovation and Infrastructure) and SDG 16 (Peace, Justice and Strong Institutions) through practical, transparent, and technology-enabled institutional operations.

Key feature contributions:
\begin{itemize}
	\item Improved access to reliable visit records for dispute resolution.
	\item Reduced manual effort through automated tracking and reports.
	\item Better institutional workflow through searchable and structured data.
\end{itemize}

\subsection{Evaluation Results and SDG Impact}
System-level testing indicates stable workflow execution for registration, detection, and logging in controlled environments. These capabilities directly support accountability and efficient administration, which in turn strengthen institutional trust and service quality.

\subsection{Future Improvements and SDG Impact}
Future upgrades such as multi-camera support, liveness detection, and centralized dashboards can broaden operational coverage and improve long-term reliability. Wider adoption across service points can further increase social and administrative benefits by reducing disputes and enabling data-driven decisions.

\chapter{Conclusion}
This mini project successfully designed and implemented a modular face detection and recognition system for college canteen visit tracking. The final prototype demonstrates how computer vision and local data systems can be integrated into a practical campus application with measurable administrative value.

The project meets its core objectives: automated identity-aware logging, support for unknown detections, configurable camera input, and accessible reporting through the desktop interface. The architecture is extensible and suitable for further academic and practical enhancement.

Future scope includes multi-camera orchestration, anti-spoofing/liveness verification, role-based access controls, cloud synchronization, and advanced analytics dashboards for long-term trend analysis.

\section*{Limitations and Future Work Breakdown}
\textbf{Current Limitations}
\begin{itemize}
	\item Performance sensitivity to poor lighting and motion blur.
	\item Single-site database deployment without distributed sync.
	\item No formal liveness/spoof detection in this iteration.
	\item Recognition thresholds require environment-specific calibration.
\end{itemize}

\textbf{Future Work Plan}
\begin{enumerate}
	\item Add liveness validation to reduce spoofing risk.
	\item Introduce multi-camera orchestration for larger coverage.
	\item Add role-based access and audit trail controls.
	\item Integrate long-term trend dashboards and anomaly alerts.
	\item Explore secure centralized storage for multi-location institutions.
\end{enumerate}

\renewcommand\bibname{References}
\addtocontents{toc}{\protect\setcounter{page}{0}}
{\thispagestyle{empty}}
\pagenumbering{gobble}
\bibliography{sample}
\thispagestyle{empty}
\addcontentsline{toc}{chapter}{References}

\newpage
\pagenumbering{roman}
\begin{appendices}

\chapter{CODE}
This appendix includes essential and representative code snippets from key modules used in the project. These snippets are intentionally expanded for documentation completeness and implementation traceability.

\section*{A.0 Module Map and Purpose}
\begin{table}[h!]
\centering
\caption{Module-level summary used in implementation}
\begin{tabular}{|l|l|}
\hline
Module & Responsibility \\
\hline
config.py & Runtime configuration, thresholds, camera/RTSP parameters \\
\hline
database.py & Schema creation, CRUD operations, log queries, daily statistics \\
\hline
face\_recognition\_module.py & Detection, enhancement, embedding extraction, recognition orchestration \\
\hline
face\_matcher.py & Similarity metrics and margin-aware best match selection \\
\hline
main.py & Console workflow, camera loop, control handlers \\
\hline
gui\_app.py & Desktop UI with tabs for detection, students, logs, and statistics \\
\hline
utils.py & Text report generation and CSV export helpers \\
\hline
test\_system.py & Interactive functional testing helper script \\
\hline
test\_rtsp.py & RTSP and webcam connectivity validation script \\
\hline
\end{tabular}
\end{table}

\section*{A.1 Configuration Snippet (config.py)}
\begin{lstlisting}[language=Python,caption={Core configuration for camera, RTSP, and recognition thresholds}]
# Database settings
DATABASE_PATH = os.path.join(BASE_DIR, "database", "canteen.db")
FACES_DIR = os.path.join(BASE_DIR, "database", "faces")
SCREENSHOTS_DIR = os.path.join(BASE_DIR, "screenshots")

# Recognition settings
FACE_RECOGNITION_THRESHOLD = 0.35
FACE_EMBEDDING_MODEL = "Facenet"
CONFIDENCE_THRESHOLD = 0.4

# Performance settings
PROCESS_EVERY_N_FRAMES = 5
DETECTION_SCALE = 1.0
USE_GPU = True
USE_THREADED_RECOGNITION = True
USE_DNN_DETECTOR = True
DNN_CONFIDENCE_THRESHOLD = 0.6

# Display and logging
SHOW_WINDOW = False
LOG_RECOGNITIONS = True
VERBOSE_RTSP_LOGS = False

# Frame freshness controls
FRAME_BUFFER_SIZE = 2
FRAME_FETCH_TIMEOUT = 0.5

# Camera settings
CAMERA_INDEX = 0
FRAME_WIDTH = 640
FRAME_HEIGHT = 480
FPS = 60

# RTSP settings
USE_RTSP = False
RTSP_URL = "rtsp://localhost:8554/live"
RTSP_RECONNECT_ATTEMPTS = 5
RTSP_RECONNECT_DELAY = 2
RTSP_BUFFER_SIZE = 100
RTSP_TRANSPORT = "tcp"
RTSP_TARGET_FPS = 7
\end{lstlisting}

\section*{A.2 Database Schema Snippet (database.py)}
\begin{lstlisting}[language=SQL,caption={Primary tables used for student management and visit logging}]
CREATE TABLE IF NOT EXISTS students (
	id INTEGER PRIMARY KEY AUTOINCREMENT,
	student_id TEXT UNIQUE NOT NULL,
	name TEXT NOT NULL,
	department TEXT,
	year INTEGER,
	face_embedding TEXT,
	face_image_path TEXT,
	created_at TIMESTAMP DEFAULT CURRENT_TIMESTAMP
);

CREATE TABLE IF NOT EXISTS visit_logs (
	id INTEGER PRIMARY KEY AUTOINCREMENT,
	student_db_id INTEGER,
	student_id TEXT,
	student_name TEXT,
	entry_time TIMESTAMP DEFAULT CURRENT_TIMESTAMP,
	is_known INTEGER DEFAULT 1
);
CREATE TABLE IF NOT EXISTS unknown_faces (
	id INTEGER PRIMARY KEY AUTOINCREMENT,
	face_image_path TEXT,
	face_embedding TEXT,
	first_seen TIMESTAMP DEFAULT CURRENT_TIMESTAMP,
	last_seen TIMESTAMP DEFAULT CURRENT_TIMESTAMP,
	times_seen INTEGER DEFAULT 1
);
\end{lstlisting}

\section*{A.3 Face Matcher Snippet (face\_matcher.py)}
\begin{lstlisting}[language=Python,caption={Similarity-based matching with confidence margin check}]
def find_best_match(embedding, known_faces, threshold=0.5, margin=0.15):
	scores = []
	for student in known_faces:
		emb = student.get('face_embedding')
		if emb is not None:
			score = cosine_similarity(embedding, emb)
			scores.append((score, student))
	if not scores:
		return None
	scores.sort(reverse=True, key=lambda x: x[0])
	best_score, best_student = scores[0]
	next_score = scores[1][0] if len(scores) > 1 else 0.0
	if best_score >= threshold and (best_score - next_score) > margin:
		return best_student
	return None
\end{lstlisting}

\section*{A.4 Main Menu Control (main.py)}
\begin{lstlisting}[language=Python,caption={Console menu and operation flow}]
while True:
	print("1. Start Real-time Detection")
	print("2. Register New Student")
	print("3. View Statistics")
	print("4. View Today's Logs")
	print("5. List All Students")
	print("6. Exit")

	choice = input("Enter your choice (1-6): ").strip()
	if choice == '1':
		app.run_detection()
	elif choice == '2':
		app.register_student_interactive()
	elif choice == '6':
		app.stop_camera()
		break
\end{lstlisting}

\section*{A.5 Core Detection Loop (main.py)}
\begin{lstlisting}[language=Python,caption={Main real-time loop with RTSP reconnect handling}]
while self.is_running:
	frame = self._get_latest_frame()
	if frame is None:
		failed_reads += 1
		if config.USE_RTSP and failed_reads >= max_failed_reads:
			self.stop_camera()
			time.sleep(2)
			if not self.start_camera():
				break
			failed_reads = 0
			continue
		time.sleep(0.1)
		continue

	failed_reads = 0
	annotated_frame, recognized_people = self.face_system.process_frame(frame)

	frame_count += 1
	elapsed_time = time.time() - start_time
	if elapsed_time >= 1.0:
		fps = frame_count / elapsed_time
		frame_count = 0
		start_time = time.time()

	if config.SHOW_WINDOW:
		cv2.imshow(config.WINDOW_TITLE, annotated_frame)
\end{lstlisting}

\section*{A.6 Face Registration from Image (face\_recognition\_module.py)}
\begin{lstlisting}[language=Python,caption={Registration flow for uploaded student image}]
def register_from_image(self, image_path, student_id, name, department, year):
	image = cv2.imread(image_path)
	if image is None:
		return False

	faces = self.detect_faces(image)
	if not faces:
		return False

	x1, y1, x2, y2 = faces[0]
	face_crop = image[y1:y2, x1:x2]
	face_enhanced = self.enhance_face(face_crop)
	embedding = self.get_face_embedding(face_enhanced)
	if embedding is None:
		return False

	face_filename = f"{student_id}_{datetime.now().strftime('%Y%m%d_%H%M%S')}_upload.jpg"
	face_path = os.path.join(config.FACES_DIR, face_filename)
	cv2.imwrite(face_path, face_enhanced, [cv2.IMWRITE_JPEG_QUALITY, 95])

	existing = database.get_student_by_id(student_id)
	if existing:
		success = database.update_student(
			student_id=student_id, name=name, department=department,
			year=year, face_embedding=embedding, face_image_path=face_path
		)
	else:
		success = database.add_student(
			student_id=student_id, name=name, department=department,
			year=year, face_embedding=embedding, face_image_path=face_path
		)
	if success:
		self.reload_known_faces()
	return success
\end{lstlisting}

\section*{A.7 GUI Student Management Snippet (gui\_app.py)}
\begin{lstlisting}[language=Python,caption={Desktop student management actions and table setup}]
def build_students_tab(self):
	toolbar = ttk.Frame(self.students_tab)
	toolbar.pack(fill=tk.X, pady=(0, 10))

	add_btn = ttk.Button(toolbar, text="Add Student", command=self.add_student_dialog)
	add_btn.pack(side=tk.LEFT, padx=5)

	upload_btn = ttk.Button(toolbar, text="Register from Image", command=self.register_from_file)
	upload_btn.pack(side=tk.LEFT, padx=5)

	refresh_btn = ttk.Button(toolbar, text="Refresh", command=self.refresh_students)
	refresh_btn.pack(side=tk.LEFT, padx=5)

	columns = ('ID', 'Student ID', 'Name', 'Department', 'Year', 'Registered')
	self.students_tree = ttk.Treeview(table_frame, columns=columns, show='headings')
\end{lstlisting}

\section*{A.11 Face Detection Strategy (face\_recognition\_module.py)}
\begin{lstlisting}[language=Python,caption={DNN-first face detection with Haar fallback}]
def detect_faces(self, frame):
	h, w = frame.shape[:2]
	result = []

	use_dnn = getattr(config, 'USE_DNN_DETECTOR', False)
	if use_dnn and self.dnn_net is not None:
		blob = cv2.dnn.blobFromImage(frame, 1.0, (300, 300), (104.0, 177.0, 123.0))
		self.dnn_net.setInput(blob)
		detections = self.dnn_net.forward()
		confidence_threshold = getattr(config, 'DNN_CONFIDENCE_THRESHOLD', 0.6)

		for i in range(detections.shape[2]):
			confidence = detections[0, 0, i, 2]
			if confidence > confidence_threshold:
				box = detections[0, 0, i, 3:7] * np.array([w, h, w, h])
				x1, y1, x2, y2 = box.astype(int)
				width = x2 - x1
				height = y2 - y1
				size = max(width, height)
				center_x = x1 + width // 2
				center_y = y1 + height // 2
				x1_sq = max(0, center_x - size // 2)
				y1_sq = max(0, center_y - size // 2)
				x2_sq = min(w, x1_sq + size)
				y2_sq = min(h, y1_sq + size)
				result.append((x1_sq, y1_sq, x2_sq, y2_sq))

		if result:
			return result

	small_frame = cv2.resize(frame, None, fx=0.5, fy=0.5, interpolation=cv2.INTER_LINEAR)
	gray = cv2.cvtColor(small_frame, cv2.COLOR_BGR2GRAY)
	faces = self.face_cascade.detectMultiScale(
		gray, scaleFactor=1.15, minNeighbors=4, minSize=(30, 30),
		flags=cv2.CASCADE_SCALE_IMAGE
	)
	for (x, y, fw, fh) in faces:
		x1 = int(x / 0.5)
		y1 = int(y / 0.5)
		result.append((x1, y1, x1 + int(fw / 0.5), y1 + int(fh / 0.5)))
	return result
\end{lstlisting}

\section*{A.12 Database Operations (database.py)}
\begin{lstlisting}[language=Python,caption={Essential CRUD and visit logging operations}]
def add_student(student_id, name, department, year, face_embedding, face_image_path):
	conn = get_connection()
	cursor = conn.cursor()
	embedding_json = json.dumps(face_embedding.tolist())
	cursor.execute('''
		INSERT INTO students (student_id, name, department, year, face_embedding, face_image_path)
		VALUES (?, ?, ?, ?, ?, ?)
	''', (student_id, name, department, year, embedding_json, face_image_path))
	conn.commit()
	conn.close()
	return True

def get_all_students():
	conn = get_connection()
	cursor = conn.cursor()
	cursor.execute('SELECT * FROM students ORDER BY name')
	rows = cursor.fetchall()
	students = []
	for row in rows:
		student = dict(row)
		if student['face_embedding']:
			student['face_embedding'] = np.array(json.loads(student['face_embedding']))
		students.append(student)
	conn.close()
	return students

def log_visit(student_db_id, student_id, student_name, screenshot_path=None, is_known=True):
	conn = get_connection()
	cursor = conn.cursor()
	cursor.execute('''
		INSERT INTO visit_logs (student_db_id, student_id, student_name, is_known, screenshot_path)
		VALUES (?, ?, ?, ?, ?)
	''', (student_db_id, student_id, student_name, 1 if is_known else 0, screenshot_path))
	log_id = cursor.lastrowid
	conn.commit()
	conn.close()
	return log_id
\end{lstlisting}

\section*{A.13 Frame Buffer and Capture Thread (main.py)}
\begin{lstlisting}[language=Python,caption={Fresh-frame strategy for reducing overlap and stale reads}]
def _start_capture_thread(self):
	self._stop_capture_thread()
	self.capture_thread_running = True

	def capture_loop():
		while self.capture_thread_running and self.cap:
			try:
				if config.USE_RTSP and config.RTSP_PREGRAB_COUNT > 0:
					for _ in range(config.RTSP_PREGRAB_COUNT):
						self.cap.grab()

				ret, frame = self.cap.read()
				if not ret:
					time.sleep(0.01)
					continue

				if self.frame_buffer.full():
					try:
						self.frame_buffer.get_nowait()
					except queue.Empty:
						pass

				try:
					self.frame_buffer.put_nowait(frame)
				except queue.Full:
					pass
			except Exception:
				time.sleep(0.01)

	self.capture_thread = Thread(target=capture_loop, daemon=True)
	self.capture_thread.start()
\end{lstlisting}

\section*{A.14 Face Enhancement Pipeline (face\_recognition\_module.py)}
\begin{lstlisting}[language=Python,caption={Image enhancement before embedding extraction}]
def enhance_face(self, face_image):
	if face_image is None or face_image.size == 0:
		return face_image

	target_size = 224
	face_resized = cv2.resize(face_image, (target_size, target_size),
							  interpolation=cv2.INTER_LANCZOS4)

	face_pil = Image.fromarray(cv2.cvtColor(face_resized, cv2.COLOR_BGR2RGB))
	enhancer = ImageEnhance.Sharpness(face_pil)
	face_pil = enhancer.enhance(1.5)
	enhancer = ImageEnhance.Contrast(face_pil)
	face_pil = enhancer.enhance(1.2)
	enhancer = ImageEnhance.Brightness(face_pil)
	face_pil = enhancer.enhance(1.1)

	face_enhanced = cv2.cvtColor(np.array(face_pil), cv2.COLOR_RGB2BGR)
	lab = cv2.cvtColor(face_enhanced, cv2.COLOR_BGR2LAB)
	l, a, b = cv2.split(lab)
	clahe = cv2.createCLAHE(clipLimit=2.0, tileGridSize=(8, 8))
	l = clahe.apply(l)
	enhanced = cv2.merge([l, a, b])
	face_enhanced = cv2.cvtColor(enhanced, cv2.COLOR_LAB2BGR)
	face_enhanced = cv2.fastNlMeansDenoisingColored(face_enhanced, None, 10, 10, 7, 21)
	return face_enhanced
\end{lstlisting}

\section*{A.15 Statistics Retrieval Logic (database.py)}
\begin{lstlisting}[language=Python,caption={Daily statistics aggregation queries}]
def get_daily_statistics(date=None):
	if date is None:
		date = datetime.now().strftime('%Y-%m-%d')

	conn = get_connection()
	cursor = conn.cursor()

	cursor.execute('SELECT COUNT(*) FROM visit_logs WHERE date = ?', (date,))
	total_visits = cursor.fetchone()[0]

	cursor.execute('SELECT COUNT(DISTINCT student_id) FROM visit_logs WHERE date = ? AND is_known = 1', (date,))
	unique_visitors = cursor.fetchone()[0]

	cursor.execute('SELECT COUNT(*) FROM visit_logs WHERE date = ? AND is_known = 0', (date,))
	unknown_visitors = cursor.fetchone()[0]

	cursor.execute('SELECT AVG(duration_minutes) FROM visit_logs WHERE date = ? AND duration_minutes IS NOT NULL', (date,))
	avg_duration = cursor.fetchone()[0] or 0

	conn.close()
	return {
		'date': date,
		'total_visits': total_visits,
		'unique_visitors': unique_visitors,
		'unknown_visitors': unknown_visitors,
		'average_duration_minutes': round(avg_duration, 2)
	}
\end{lstlisting}

\section*{A.16 GUI Tab Construction Overview (gui\_app.py)}
\begin{lstlisting}[language=Python,caption={Notebook tabs used in desktop interface}]
def build_ui(self):
	main_container = ttk.Frame(self.root, padding="10")
	main_container.pack(fill=tk.BOTH, expand=True)

	self.notebook = ttk.Notebook(main_container)
	self.notebook.pack(fill=tk.BOTH, expand=True)

	self.detection_tab = ttk.Frame(self.notebook)
	self.notebook.add(self.detection_tab, text="Live Detection")
	self.build_detection_tab()

	self.students_tab = ttk.Frame(self.notebook)
	self.notebook.add(self.students_tab, text="Students")
	self.build_students_tab()

	self.logs_tab = ttk.Frame(self.notebook)
	self.notebook.add(self.logs_tab, text="Visit Logs")
	self.build_logs_tab()

	self.stats_tab = ttk.Frame(self.notebook)
	self.notebook.add(self.stats_tab, text="Statistics")
	self.build_stats_tab()
\end{lstlisting}

\section*{A.17 Export Utilities (utils.py)}
\begin{lstlisting}[language=Python,caption={CSV export logic for logs and students}]
def export_logs_to_csv(filepath, date=None, student_id=None):
	logs = database.get_visit_logs(date=date, student_id=student_id)
	with open(filepath, 'w', newline='', encoding='utf-8') as f:
		f.write("ID,Date,Entry Time,Exit Time,Student ID,Student Name,Duration (min),Status\n")
		for log in logs:
			status = "Known" if log.get('is_known') else "Unknown"
			duration = log.get('duration_minutes', '')
			f.write(f"{log['id']},{log.get('date', '')},{log.get('entry_time', '')},{log.get('exit_time', '')},")
			f.write(f"{log.get('student_id', '')},\"{log.get('student_name', '')}\",{duration},{status}\n")

def export_students_to_csv(filepath):
	students = database.get_all_students()
	with open(filepath, 'w', newline='', encoding='utf-8') as f:
		f.write("ID,Student ID,Name,Department,Year,Registered Date\n")
		for student in students:
			created = student.get('created_at', '')[:10] if student.get('created_at') else ''
			f.write(f"{student['id']},{student['student_id']},\"{student['name']}\",")
			f.write(f"{student.get('department', '')},{student.get('year', '')},{created}\n")
\end{lstlisting}

\section*{A.18 Startup Setup Automation (setup.py)}
\begin{lstlisting}[language=Python,caption={First-run setup routine}]
def install_requirements():
	packages = [
		"opencv-python>=4.8.0",
		"numpy>=1.24.0",
		"Pillow>=10.0.0",
		"ultralytics>=8.1.0",
		"deepface>=0.0.79",
		"tensorflow>=2.13.0"
	]
	for package in packages:
		try:
			subprocess.check_call([sys.executable, "-m", "pip", "install", package, "-q"])
		except subprocess.CalledProcessError:
			pass

	base_dir = os.path.dirname(os.path.abspath(__file__))
	os.makedirs(os.path.join(base_dir, "database", "faces"), exist_ok=True)
	os.makedirs(os.path.join(base_dir, "screenshots"), exist_ok=True)

	import database
	database.init_database()
\end{lstlisting}

\section*{A.19 Detection Overlay Utility (utils.py)}
\begin{lstlisting}[language=Python,caption={Bounding box and label overlay utility}]
def draw_face_box(frame, bbox, name, confidence, is_known):
	x1, y1, x2, y2 = bbox
	if is_known:
		color = (0, 255, 0)
		bg_color = (0, 200, 0)
	else:
		color = (0, 0, 255)
		bg_color = (0, 0, 200)

	cv2.rectangle(frame, (x1, y1), (x2, y2), color, 2)
	label = f"{name} ({confidence:.0%})"
	(tw, th), bl = cv2.getTextSize(label, cv2.FONT_HERSHEY_SIMPLEX, 0.6, 2)
	cv2.rectangle(frame, (x1, y1 - th - 10), (x1 + tw + 10, y1), bg_color, -1)
	cv2.putText(frame, label, (x1 + 5, y1 - 5), cv2.FONT_HERSHEY_SIMPLEX, 0.6, (255, 255, 255), 2)
	return frame
\end{lstlisting}

\section*{A.20 Additional Screenshot Spaces for Final Documentation}
\begin{figure}[h!]
	\begin{center}
		\fbox{\rule{0pt}{2.9in}\rule{0.90\linewidth}{0pt}}
		\caption{Additional Slot I: Database table browser evidence}
	\end{center}
\end{figure}

\begin{figure}[h!]
	\begin{center}
		\fbox{\rule{0pt}{2.9in}\rule{0.90\linewidth}{0pt}}
		\caption{Additional Slot J: Student list management view}
	\end{center}
\end{figure}

\begin{figure}[h!]
	\begin{center}
		\fbox{\rule{0pt}{2.9in}\rule{0.90\linewidth}{0pt}}
		\caption{Additional Slot K: Statistics tab output}
	\end{center}
\end{figure}

\begin{figure}[h!]
	\begin{center}
		\fbox{\rule{0pt}{2.9in}\rule{0.90\linewidth}{0pt}}
		\caption{Additional Slot L: CSV export file and report snapshot}
	\end{center}
\end{figure}

\section*{A.8 RTSP Connectivity Test Snippet (test\_rtsp.py)}
\begin{lstlisting}[language=Python,caption={RTSP stream test before full system execution}]
def test_rtsp_stream(rtsp_url):
    cap = cv2.VideoCapture(rtsp_url, cv2.CAP_FFMPEG)
    if not cap.isOpened():
        return False

    frame_count = 0
    start_time = time.time()
    for i in range(30):
        ret, frame = cap.read()
        if ret:
            frame_count += 1

    elapsed = time.time() - start_time
    fps = frame_count / elapsed if elapsed > 0 else 0
    cap.release()
    return frame_count >= 25
\end{lstlisting}

\section*{A.9 Utility Reporting Snippet (utils.py)}
\begin{lstlisting}[language=Python,caption={Daily report creation logic}]
def generate_report(start_date=None, end_date=None):
	if start_date is None:
		start_date = datetime.now().strftime('%Y-%m-%d')
	if end_date is None:
		end_date = start_date

	report = []
	report.append("=" * 60)
	report.append("COLLEGE CANTEEN FACE DETECTION SYSTEM - REPORT")
	report.append("=" * 60)
	report.append(f"Report Period: {start_date} to {end_date}")

	stats = database.get_daily_statistics(start_date)
	report.append(f"Total Visits: {stats['total_visits']}")
	report.append(f"Unique Visitors: {stats['unique_visitors']}")
	report.append(f"Unknown Visitors: {stats['unknown_visitors']}")
\end{lstlisting}

\section*{A.10 Testing Script Snippet (test\_system.py)}
\begin{lstlisting}[language=Python,caption={Interactive test menu for registration and recognition validation}]
while True:
	print("1. Register from IMAGE FILE")
	print("2. Register ALL from FOLDER")
	print("3. Register from WEBCAM")
	print("4. Test with WEBCAM")
	print("5. View registered students")
	print("6. Re-register (update face)")
	print("7. Exit")

	choice = input("Enter choice (1-7): ").strip()
	if choice == '4':
		test_with_webcam(system)
	elif choice == '7':
		break
\end{lstlisting}

\chapter{Screenshots}
Screenshots to be inserted in the final submission. Dedicated placeholder areas are provided below to avoid last-minute formatting shifts.

\begin{figure}[h!]
	\begin{center}
		\fbox{\rule{0pt}{2.8in}\rule{0.88\linewidth}{0pt}}
		\caption{Screenshot Slot A: Main dashboard}
	\end{center}
\end{figure}

\begin{figure}[h!]
	\begin{center}
		\fbox{\rule{0pt}{2.8in}\rule{0.88\linewidth}{0pt}}
		\caption{Screenshot Slot B: Live detection with labels}
	\end{center}
\end{figure}

\begin{figure}[h!]
	\begin{center}
		\fbox{\rule{0pt}{2.8in}\rule{0.88\linewidth}{0pt}}
		\caption{Screenshot Slot C: Student registration form}
	\end{center}
\end{figure}

\begin{figure}[h!]
	\begin{center}
		\fbox{\rule{0pt}{2.8in}\rule{0.88\linewidth}{0pt}}
		\caption{Screenshot Slot D: Visit logs and filters}
	\end{center}
\end{figure}

\begin{figure}[h!]
	\begin{center}
		\fbox{\rule{0pt}{2.8in}\rule{0.88\linewidth}{0pt}}
		\caption{Screenshot Slot E: Statistics and summary panel}
	\end{center}
\end{figure}

\begin{figure}[h!]
	\begin{center}
		\fbox{\rule{0pt}{2.8in}\rule{0.88\linewidth}{0pt}}
		\caption{Screenshot Slot F: RTSP stream capture window}
	\end{center}
\end{figure}

\begin{figure}[h!]
	\begin{center}
		\fbox{\rule{0pt}{2.8in}\rule{0.88\linewidth}{0pt}}
		\caption{Screenshot Slot G: RTSP test and reconnect logs}
	\end{center}
\end{figure}

\begin{figure}[h!]
	\begin{center}
		\fbox{\rule{0pt}{2.8in}\rule{0.88\linewidth}{0pt}}
		\caption{Screenshot Slot H: Exported CSV/report evidence}
	\end{center}
\end{figure}

Screenshots checklist:
\begin{itemize}
	\item Desktop dashboard home screen.
	\item Live detection view with recognized and unknown labels.
	\item Student registration dialog and success output.
	\item Visit logs tab with filters and exported CSV sample.
	\item Statistics tab and daily summary panel.
	\item RTSP stream capture screen and status.
	\item RTSP test output and reconnection behavior evidence.
\end{itemize}

\end{appendices}

\end{document}
